\documentclass{exam}
\usepackage{amsmath, amssymb}

\pagestyle{headandfoot}
\firstpageheadrule
\runningheadrule
\firstpageheader{Prof. Rong \\ Introduction to Real Analysis I}{Homework 6}{Aadhithya Saravanan}
\runningheader{Introduction to Real Analysis I \\ Homework 6}{}{Saravanan}
\firstpagefooter{}{}{}
\runningfooter{ }{\thepage}{ }

\printanswers

\begin{document}

\underline{Exercise 3.2.2}
\newline
\begin{parts}
    \part A point $x$ is a limit point of $A$ if for every $\epsilon>0$, $V_\epsilon(x)$ intersects $A$ at some point other than $x$. This is equivalent to there being a sequence of elements in $A$ that converge to $x$. The limit points of $A$ are $-1$ and 1. A sequence that converges to $-1$ is the set of points in $A$ for odd $n$. Let $a_n=-1+\frac{2}{n}$, a sequence of terms from $A$. As $n\to\infty$, $a_n\to{-1}$. A sequence that converges to 1 is the set of points in $A$ for even $n$. Let $b_n=1+\frac{2}{n}$, a sequence of terms from $A$. As $n\to\infty$, $b_n\to1$. Now, we need to show that no other real number is a limit point of $A$. Any subsequence of $A$ has either infinitely many terms from $a_n$ or infitely many terms from $b_n$ or contains infinitely many terms from both. $x\in\mathbb{R}$ with $x\ne -1$ and $x\ne1$ is a limit point of $A$ if and only if there exists a sequence $(x_n)$ with terms from $A\backslash\{x\}$ that converges to $x$. Assume $x$ is a limit point of $A$. Then, by Theorem 3.2.5, there exists a sequence $(x_n)$ in $A\backslash\{x\}$ that converges to $x$. $A$ is the union of the two disjoint sets $C=\{-1+\frac{2}{n}\}$ for odd $n$ and $D=\{1+\frac{2}{n}\}$ for even $n$. Choose any such $(x_n)$. It must have infinitely many terms from $C$ or $D$. For the case where $(x_n)$ has infinitely many terms from $C$, we can take the terms from $(x_n)$ that are in $C$ in the order they appear in $(x_n)$ to form a subsequence $(x_{n_C})$. Since $(x_n)$ converges to $x$, $(x_{n_C})$ must also converge to $x$. But, $(x_{n_C})$ is contained in $C$, which has only $-1$ as a limit point. This is because $a_n=-1+\frac{2}{n}$, which converges to $-1$, contains all points in $C$, so any convergent subsequence of points, such as $(x_{n_C})$, from $C$ must converge to $-1$. Since a subsequence of $(x_n)$ converges to $-1$ and $(x_n)$ converges, $(x_n)$ converges to $-1$. This contradicts the assumption that $x\ne-1$. The same argument can be used to show that if $(x_n)$ has infinitely many terms from $D$, it converges to 1, contradicting that $x\ne1$. So, in either case, the assumption that $x$ is a limit point must be false. Therefore, $-1$ and 1 are the only limit points of $A$.
    \newline
    \newline
    The limit points of $B$ are all values in the set $\{0\le x\le1|x\in\mathbb{R}\}$. Choose an arbitary $x$ from this set. By the Density of $\mathbb{Q}$ in $\mathbb{R}$, for any $\epsilon>0$, there exists a rational number $q\ne x$ in $V_\epsilon(x)=(x-\epsilon,x+\epsilon)$. So, taking the intersection of $V_\epsilon(x)=(x-\epsilon,x+\epsilon)$ and $B$ must have this $q$, which is an element of $B$ that is not equal to the arbitrarily chosen $x$. For any point $y$ outside of $\{0\le x\le1|x\in\mathbb{R}\}$, there exists an $\epsilon>0$ such that the intersection of $V_\epsilon(y)$ and $B$ is empty. For $y<0$, choose $\epsilon=-\frac{y}{2}$. Then, $V_\epsilon(y)=(\frac{3}{2}y,\frac{1}{2}y)$ is a neighborhood that only contains negative values, and therefore the intersection of $V_\epsilon(y)$ and $B$ is empty. For $y>1$, choose $\epsilon=\frac{y-1}{2}$. Then, $V_\epsilon(y)=(\frac{y+1}{2},\frac{3y-1}{2})$ is a neighborhood that only contains values greater than 1, and therefore the intersection of $V_\epsilon(y)$ and $B$ is empty. So, the only limit points of $B$ are the values in the set $\{0\le x\le1|x\in\mathbb{R}\}$.
    \part $A$ is not open nor closed. For $A$ to be open, for each element $a\in A$, there must exist a $V_\epsilon(a)$ that is a subset of $A$. For any $\epsilon>0$, we can always find a real number that is distance of less than $\epsilon$ from $a$. Since every element of $A$ is rational, any real number from $V_\epsilon(a)$ that is irrational shows that $V_\epsilon(a)$ is not a subset of $A$. Therefore, $A$ is not open. For $A$ to be closed, it must contain all of its limit points. Since $A$ does not contain $-1$ or 1, its limit points, it is not closed. $B$ is also not open nor closed. For any $\epsilon>0$, we can always find a real number that is distance of less than $\epsilon$ from any $b\in B$. Since every element of $B$ is rational, any real number from $V_\epsilon(b)$ that is irrational shows that $V_\epsilon(b)$ is not a subset of $B$. The limit points of $B$ are the values in the set $\{0\le x\le1|x\in\mathbb{R}\}$, which includes irrationals. Therefore, $B$ does not contain its limit points and is not closed.
    \part All points of $A$ are isolated. For an element $a\in A$ to be isolated, it cannot be a limit point of $A$. The only limit points of $A$ are $-1$ and 1, which are not in $A$. So, all elements of $A$ are isolated. No point of $B$ is isolated. The elements of $B$ are a subset of $\{0\le x\le1|x\in\mathbb{R}\}$, which is the set of limit points of $B$. Therefore, every point in $B$ is a limit point and there are no isolated points.
    \part The closure of $A$ is the set $\bar{A}=A\cup\{-1,1\}$. Therefore, the closure of $A$ is all elements of $A$ and the points $-1$ and 1. The closure of $B$ is the set $\bar{B}=B\cup\{0\le x\le1|x\in\mathbb{R}\}$. Since $B$ is a subset of $\{0\le x\le1|x\in\mathbb{R}\}$, the closure is simply $\{0\le x\le1|x\in\mathbb{R}\}$.
    \newline
\end{parts}

\underline{Exercise 3.2.4}
\newline
\begin{parts}
    \part By definition of supremum, for any $\epsilon>0$, $s-\epsilon<a$ for some $a\in A$. Let $\epsilon_n=\frac{1}{n}$. Choose $a_n$ as an element in $A$ satisfying $s-\epsilon_n<a_n$. This limit of this sequence $(a_n)$ is a limit point of $A$, as it is a sequence made up of terms in $A$. It converges to $s$ as $n\to\infty$. By definition of closure, $\bar{A}$ contains all limit points of $A$, so $s\in\bar{A}$.
    \part Let $A$ be an open set containing its supremum, $s$. For any $\epsilon>0$, $V_\epsilon(s)=(s-\epsilon,s+\epsilon)$. Choose any real number $x\in(s,s+\epsilon)$. This element is an element of $V_\epsilon(s)$, but $x>s>=a,\forall a\in A$. So, $x\notin A$. Since this element is in $V_\epsilon(s)$ but not in $A$, $V_\epsilon(s)\nsubseteq A$. By definition of open sets, $A$ cannot be open, so we arrive at contradiction. Therefore, open sets cannot contain their supremums.
    \newline
\end{parts}

\underline{Exercise 3.2.6}
\newline
\begin{parts}
    \part Let $O$ be $\{\mathbb{R}\backslash{\sqrt{2}}\}$. $O$ contains all rational numbers, as $\mathbb{R}$ contains all the rationals and removing an irrational number keeps all the rationals within $O$. $O^c=\{\sqrt{2}\}$. $O^c$ contains no limit points and is therefore closed. By Theorem 3.2.13, since $O^c$ is closed, $O$ is open.
    \part To prove this false, we are looking for a collection of nested closed sets $O_1\supseteq O_2\supseteq O_3\supseteq\dots$ such that $\bigcap_{n=1}^\infty O_n=\emptyset$. Let $O_n=[n,\infty)$. $O_n$ is a closed set, as the limit points for each $O_n$ are all elements in $O_n$. Then for any $x\in\mathbb{R}$, choose $m\in\mathbb{N}$ such that $m$ is the least natural number greater than $x$. Then, $x\notin[m,\infty)=O_m$. So, $x\notin\bigcap_{n=1}^\infty O_n$. Since $x$ was chosen arbitrarily from $\mathbb{R}$, there is no real number in $\bigcap_{n=1}^\infty O_n$. So, $\bigcap_{n=1}^\infty O_n=\emptyset$.
    \part Assume that there exists a nonempty open set $P$ without any rational numbers. Then, for each $a\in P$, $V_\epsilon(a)\subseteq P$ for some $\epsilon>0$. $V_\epsilon(a)=(a-\epsilon,a+\epsilon)$. By the Density of $\mathbb{Q}$ in $\mathbb{R}$, $\exists q\in\mathbb{Q}$ such that $q\in(a-\epsilon,a+\epsilon)=V_\epsilon(a)$. Since $V_\epsilon(a)\subseteq P$, $q\in P$. However, this contradicts the assumption that $P$ does not contain any rational numbers. Therefore, our assumption is wrong and every nonempty open set must contain a rational number.
    \part To prove this false, we are looking for a bounded infinite closed set whose limit point is in the set and contains no rational numbers. Let $A=\{\sqrt{2}+\frac{1}{n}|n\in\mathbb{N}\}\cup\{\sqrt{2}\}$. $\sqrt{2}+1$ is the greater number in $A$ and $\sqrt{2}$ is a lower bound for $A$, so $A$ is bounded. This set is infinite because there is a term corresponding to each natural number. A limit point of $A$ is $\sqrt{2}$. To show this, we can take the sequence $(a_n)$ where $a_n=\sqrt{2}+\frac{1}{n}$. This converges to $\sqrt{2}$ and $a_n\ne\sqrt{2}$ for all $n$, so by Theorem 3.2.5, $\sqrt{2}$ is a limit point of $A$. To show that $\sqrt{2}$ is the only limit point of $A$, we can use a proof by contradiction. Assume $x$ is a limit point of $A$ and $x\ne\sqrt{2}$. Then, by Theorem 3.2.5, there exists a sequence $(x_n)$ contained in $A$ with $x_n\ne x$ and $(x_n)\to x$. $(x_n)$ contains an infinite number of terms from $\{\sqrt{2}+\frac{1}{n}|n\in\mathbb{N}\}$. Using those terms, we have a subsequence $(x_{n_k})$, which must converge to the same limit as $(x_n)$, $x$. Since $(x_{n_k})$ is contained in $\{\sqrt{2}+\frac{1}{n}|n\in\mathbb{N}\}$, which converges to $\sqrt{2}$. So, $x=\sqrt{2}$, contradicting the assumption that $x\ne\sqrt{2}$. So, $\sqrt{2}$ is the only limit point of $A$. $A$ has no rational numbers, as it is formed by an irrational number and the sum of irrational numbers and rational numbers, which are irrational numbers. So, bounded infinite sets do not necessarily contain a rational number.
    \part The Cantor set is closed. $C_0=[0,1]$ is a closed set. $C_1$ is constructed by removing the middle third of $C_0$, resulting in $C_1=[0,\frac{1}{3}]\cup[\frac{2}{3},1]$. This set is closed by Theorem 3.2.14, as it is a union of a finite collection of closed sets. This continues for each $C_n$, where each $C_n$ is closed because it is a union of a finite collection of closed sets. Since the Cantor set is the intersection of all $C_n$, it is closed by Theorem 3.2.14, as it is the intersection of an arbitrary collection of closed sets.
    \newline
\end{parts}

\underline{Exercise 3.2.9}
\newline
\begin{parts}
    \part First, I will show that $(\bigcup_{\lambda\in\Lambda}E_\lambda)^c=\bigcap_{\lambda\in\Lambda}E_\lambda^c$. Let $x\in(\bigcup_{\lambda\in\Lambda}E_\lambda)^c$. Then, by definition of complements, $x\notin\bigcup_{\lambda\in\Lambda}E_\lambda$. By the definition of an arbitrary union, $x\notin E_\lambda$ for all $\lambda\in\Lambda$. So, by the definition of complements, $x\in E_\lambda^c$ for each $\lambda\in\Lambda$. Since $x$ is in every $E_\lambda^c$, it is in their intersection by the definition of arbitrary intersections, so $x\in\bigcap_{\lambda\in\Lambda}E_\lambda^c$. Since each step applies in both directions, $(\bigcup_{\lambda\in\Lambda}E_\lambda)^c=\bigcap_{\lambda\in\Lambda}E_\lambda^c$.
    \newline
    \newline
    Now, I will show that $(\bigcap_{\lambda\in\Lambda}E_\lambda)^c=\bigcup_{\lambda\in\Lambda}E_\lambda^c$. Let $x\in(\bigcap_{\lambda\in\Lambda}E_\lambda)^c$. Then, by definition of complements, $x\notin\bigcap_{\lambda\in\Lambda}E_\lambda$. By the definition of an arbitrary intersection, $x\notin E_\lambda$ for at least one $\lambda_0\in\Lambda$. So, by the definition of complements, $x\in E_{\lambda_0}^c$. Since $x$ is in $E_{\lambda_0}^c$ for $\lambda_0\in\Lambda$, $x\in\bigcup_{\lambda\in\Lambda}E_\lambda^c$ by the definition of arbitrary unions. Since each step applies in both directions, $(\bigcap_{\lambda\in\Lambda}E_\lambda)^c=\bigcup_{\lambda\in\Lambda}E_\lambda^c$.
    Hence, De Morgan's laws apply to any number of sets.
    \part First, I will show that the union of a finite collection of closed sets is closed. Let $\{O_1,O_2,O_3,\cdots O_N\}$ be a finite collection of closed sets. We want to show that $O=\bigcup_{n=1}^NO_n$ is closed. By Theorem 3.2.13, this is equivalent to showing that $O^c=(\bigcup_{n=1}^NO_n)^c$ is open. With De Morgan's Laws, $O^c=\bigcap_{n=1}^NO_n^c$. Since each $O_n$ is closed, each $O_n^c$ is open by Theorem 3.2.13. Since $O^c$ is the intersection of a finite collection of open sets, $O^c$ is open by Theorem 3.2.3. Therefore, the union of a finite collection of closed sets is closed.
    \newline
    \newline
    Now, I will show that the intersection of an arbitrary collection of closed sets is closed. Let $\{O_\lambda|\lambda\in\Lambda\}$ be an arbitrary collection of closed sets. We want to show that $O=\bigcap_{\lambda\in\Lambda}O_\lambda$ is closed. By Theorem 3.2.13, this is equivalent to showing that $O^c=(\bigcap_{\lambda\in\Lambda}O_\lambda)^c$ is open. With De Morgan's Laws, $O^c=\bigcup_{\lambda\in\Lambda}O_\lambda^c$. Since each $O_\lambda$ is closed, each $O_\lambda^c$ is open by Theorem 3.2.13. Since $O^c$ is the union of an arbitrary collection of open sets, $O^c$ is open by Theorem 3.2.3. Therefore, the intersection of an arbitrary collection of closed sets is closed.
    \newline
\end{parts}

\underline{Exercise 3.2.15}
\newline
\begin{parts}
    \part To show that $[a,b]$ is a $G_\delta$ set, we must represent $[a,b]$ as a countable intersection of open sets. Define the open intervals $A_n=(a-\frac{1}{n},b+\frac{1}{n})$ for $n\in \mathbb{N}$. For each $n\in \mathbb{N}$, we have a corresponding $A_n$, so the number of sets $A_n$ is countable. From Example 3.2.2.ii, any open interval on $\mathbb{R}$ is open, so each $A_n$ is open. Let $A$ be the intersection of all $A_n$. To show that $A=[a,b]$, I will first show that $a\in A$. For each $A_n=(a-\frac{1}{n},b+\frac{1}{n})$, $a>a-\frac{1}{n}$, as $\frac{1}{n}>0$. So, $a\in A_n$ for all $A_n$. Therefore, $a\in\bigcap_{n\in \mathbb{N}} A_n=A$. Now, I will show that $b\in A$. For each $A_n=(a-\frac{1}{n},b+\frac{1}{n})$, $b<b+\frac{1}{n}$, as $\frac{1}{n}>0$. So, $b\in A_n$ for all $A_n$. Therefore, $b\in\bigcap_{n\in \mathbb{N}} A_n=A$. Any $x\in \mathbb{R}$ with $a<x<b$ is clearly in $A$, as $a-\frac{1}{n}<a<x<b<b+\frac{1}{n}$, so $x\in A_n$ for all $A_n$. To show that no value less than $a$ is in $A$, pick any $y\in \mathbb{R}$ with $y<a$. Choose any $n\in \mathbb{N}$ with $n>\frac{1}{a-y}$. Then, we can rearrange to see that $y<a-\frac{1}{n}$, so $y\notin A_n$ for this choice of $n$. To show that no value greater than $b$ is in $A$, pick any $z\in \mathbb{R}$ with $z>b$. Choose any $n\in \mathbb{N}$ with $n>\frac{1}{z-b}$. Then, we can rearrange to see that $z>b+\frac{1}{n}$, so $z\notin A_n$ for this choice of $n$. So, $z\notin\bigcap_{n\in \mathbb{N}} A_n=A$. Therefore, $A=[a,b]$. Since $[a,b]$ is the intersection of all $A_n$, $[a,b]$ is the countable intersection of open sets and is a $G_\delta$ set.
    \part First, we will show that $(a,b]$ is a $G_\delta$ set, so we must represent $(a,b]$ as a countable intersection of open sets. Let $A_n=(a,b+\frac{1}{n})$. For each $n\in \mathbb{N}$, we have a corresponding $A_n$, so the number of sets $A_n$ is countable. From Example 3.2.2.ii, any open interval on $\mathbb{R}$ is open, so each $A_n$ is open. Let $A$ be the intersection of all $A_n$. $a\notin A_n$, as $A_n$ is an open interval, so $a\notin A$. Clearly, no $x<a$ is in any $A_n$, so $x\notin A$. Now, I will show that $b\in A$. For each $A_n=(a,b+\frac{1}{n})$, $b<b+\frac{1}{n}$, as $\frac{1}{n}>0$. So, $b\in A_n$ for all $A_n$. Therefore, $b\in\bigcap_{n\in \mathbb{N}} A_n=A$. Any $y\in \mathbb{R}$ with $a<y<b$ is clearly in $A$, as $a<y<b<b+\frac{1}{n}$, so $y\in A_n$ for all $A_n$. To show that no value greater than $b$ is in $A$, pick any $y\in \mathbb{R}$ with $y>b$. Choose any $n\in \mathbb{N}$ with $n>\frac{1}{y-b}$. Then, we can rearrange to see that $y>b+\frac{1}{n}$, so $y\notin A_n$ for this choice of $n$. So, $y\notin\bigcap_{n\in \mathbb{N}} A_n=A$. Therefore, $A=(a,b]$. Since $(a,b]$ is the intersection of all $A_n$, $(a,b]$ is the countable intersection of open sets and is a $G_\delta$ set.
    \newline
    \newline
    Now, I will show that $(a,b]$ is a $F_\sigma$ set, so we must represent $(a,b]$ as a countable union of closed sets. Let $B_n=[a+\frac{1}{n},b]$. For each $n\in \mathbb{N}$, we have a corresponding $B_n$, so the number of sets $B_n$ is countable. From Example 3.2.9.ii, any closed interval on $\mathbb{R}$ is closed, so each $B_n$ is closed. Let $B$ be the union of all $B_n$. $a\notin B_n$ for all $B_n$, as $a+\frac{1}{n}>a$ because $\frac{1}{n}>0$. So, $a\notin B$. Clearly, $b$ is in $B$, as $b\in B_n$ for any $B_n$. Any $x\in \mathbb{R}$ with $a<x<b$ is in $B$. Choose any $n\in \mathbb{N}$ with $n>\frac{1}{x-a}$. Then, we can rearrange to see that $x>a+\frac{1}{n}$. So, $a+\frac{1}{n}<x<b$ and $x\in B_n$. Therefore, $x\in B$. Clearly, no value greater than $b$ is in $B$, as each $B_n$ has $b$ as an upper bound. Also, no value less than $a$ is in $B$, as each $B_n$ has $a$ as a lower bound. Since $(a,b]$ is the union of all $B_n$, $(a,b]$ is the countable union of closed sets and is a $F_\sigma$ set.
    \part Every singleton set among the reals is closed. This is because the complement of a singleton set contains every real number except for one, making it an open set. We also know that there is a bijection from $\mathbb{Q}$ to $\mathbb{N}$, so the number of rational numbers is countable. Take the collection of all singleton sets with only a rational number. This collection is countable because there are only a countable number of rational numbers. Each set is also closed because it is a singleton set. Taking the union of all sets in this collection, we get $\mathbb{Q}$. So, $\mathbb{Q}$ can be written as the countable union of closed sets.
    \newline
    \newline
    Using the same collection of singleton sets as in the last part, we can take the complement of each set consisting of exactly one rational number. Since each set was closed, the complement of each set is open. Looking at the collection of all such sets, for each $q\in\mathbb{Q}$, there exists a set $O_q$ in the collection such that $O_q=\mathbb{R}\backslash\{q\}$. The number of sets in this collection is countable, as there is a bijective mapping from each rational number to a set in the collection, and there are a countable number of rational numbers. Taking the intersection of all the sets in this collection, we have $\mathbb{R}\backslash\mathbb{Q}=\mathbb{I}$. This is because for each $q\in\mathbb{Q}$, there exists a set in the collection with $O_q=\mathbb{R}\backslash\{q\}$, so an intersection with this set does not contain $q$. Therefore, $\mathbb{I}$ can be written as the countable intersection of open sets.
    \newline
\end{parts}

\underline{Exercise 3.3.1}
\newline
\newline
By Theorem 3.3.5, if $K$ is compact, it is closed and bounded. Since $K$ is bounded, there exists an $M\in \mathbb{R}$ such that $|k|<M$ for all $k\in K$. Since $K$ is nonempty and is bounded both above and below, $s=\sup K$ and $i=\inf K$ exist. By the definition of supremum, for any $\epsilon>0$, $s-\epsilon<k\le s$ for some $k\in K$. Let $\epsilon_n=\frac{1}{n}>0$ for $n\in \mathbb{N}$ and $k_n$ be a corresponding point in $K$ satisfying $s-\epsilon_n<k_n\le s$. Then, $(k_n)$ is a sequence whose limit is $s$. So, $s$ is a limit point of $K$. Since $K$ is closed, $s\in K$. A similar argument can be used to show that $i\in K$, using the definition of infimum to have $i+\epsilon>k\ge i$ for some $k\in K$. So, $\sup K$ and $\inf K$ exist and are in $K$.
\newline

\underline{Exercise 3.3.2}
\newline
\begin{parts}
    \part $\mathbb{N}$ is not compact because it is not bounded. An example of a sequence in $\mathbb{N}$ that has no convergent subsequence is $(a_n)$ where $a_n=n$ for $n\in \mathbb{N}$. The sequence is increasing and unbounded, so any subsequence is also increasing and unbounded. Therefore, $(a_n)$ has no subsequence that converges to a limit in $\mathbb{N}$, so $\mathbb{N}$ is not compact.
    \part $\mathbb{Q}\cap[0,1]$ is not compact. By Theorem 3.2.10, there exists a sequence $(a_n)$ with terms from $\mathbb{Q}$ that converges to $\frac{\sqrt{2}}{2}\in \mathbb{R}$. Let $(a_{n_k})$ be the subsequence of this sequence with terms in $\mathbb{Q}\cap[0,1]$. $(a_{n_k})$ is convergent to the same limit as $(a_n)$, which is $\frac{\sqrt{2}}{2}$. Since $(a_{n_k})$ is convergent, all of its subsequences converge to $\frac{\sqrt{2}}{2}\notin\mathbb{Q}\cap[0,1]$ and $\mathbb{Q}\cap[0,1]$ is not compact.
    \part The Cantor set is compact. Since it is closed and bounded by 1, by Theorem 3.3.4, it is compact.
    \part Let $(a_n)$ be the sequence of terms containing all values in the set $\{1+\frac{1}{2^2}+\frac{1}{3^2}+\dots+\frac{1}{n^2}|n\in \mathbb{N}\}$ with $a_n=1+\frac{1}{2^2}+\frac{1}{3^2}+\dots+\frac{1}{n^2}$. Each $a_n$ is rational because every term in the set is a rational, as each term is rational, being a finite sum of rational numbers. However, the infinite sum is known to be convergent to an irrational number. Since $(a_n)$ is convergent, any subsequence of $(a_n)$ converges to the same irrational number and has a limit not in the set. So, $\{1+\frac{1}{2^2}+\frac{1}{3^2}+\dots+\frac{1}{n^2}|n\in \mathbb{N}\}$ is not a compact set.
    \part The set is compact, since every convergent subsequence of all sequences made up of terms from the set converges to 1.
    \newline
\end{parts}

\underline{Exercise 3.3.5}
\newline
\begin{parts}
    \part Since compact sets are closed by Theorem 3.3.4, we can use Theorem 3.2.14, to say that the arbitrary intersection of compact sets is closed. The intersection is a subset of each of the compact sets, so the intersection is bounded. Therefore, the arbitrary intersection of compact sets is compact.
    \part Let $A_n=[n,n+1]$ for $n\in \mathbb{N}$. Each $A_n$ is compact, as Example 3.3.2 showed that any closed interval is a compact set. Taking the union of all $A_n$, we have $\bigcup_{n=1}^\infty A_n=[1,\infty)$. Any positive real number $x>1$ is in the infinite union, as $A_n$ with $n=\lfloor x\rfloor$ contains $x$. So, the union is not bounded, and by Theorem 3.3.4, the union is not compact.
    \part Let $A=(0,1)$ and $K=[0,1]$. Then, $A\cap K=(0,1)$ is a set that is not closed, as it has limit points of 0 and 1 that are not in the set. By Theorem 3.3.4, the intersection is not compact.
    \part Let $F_n=[n,\infty)$ for $n\in \mathbb{N}$. Then, $\bigcap_{n=1}^\infty F_n=\emptyset$. For any real number $x$, choose any $n\in \mathbb{N}$ with $n>x$. Then, $F_n$ does not contain $x$, so $x\notin\bigcap_{n=1}^\infty F_n$. So, $\bigcap_{n=1}^\infty F_n=\emptyset$.
    \newline
\end{parts}

\end{document}